\documentclass[exercice,a4]{thedocument}%
\usepackage[french]{babel}%
\usepackage{tabularx}%
\usepackage{array}%
\usepackage[np]{numprint}%
\startdatakeys
    \source{22GENMATAA1}
    \BO{2026}
    \cycle{4}
    \competencesDuSocle{CA,RA}
    \objectifsApprentissage{C4C2ab,C4D2ae,C4C2ag}
    \motsCles{tableur,pourcentage,pourcentage d'augmentation}
\enddatakeys
\begin{document}%

\medskip
Lors des Jeux paralympiques de 2021, les médias ont proposé un classement des pays en
fonction de la répartition des médailles obtenues. Voici le classement obtenu pour les 15 premiers pays :

\begin{center}
\begin{tabularx}{\linewidth}{|c|m{2.75cm}|*{5}{>{\centering \arraybackslash}X|}}\hline
&A  &B  &C  &D  &E  &F\\ \hline
	1  & Nations &Classement &Or &Argent &Bronze &Total\\ \hline
	2  & Chine      &1  &96 &60 &51 &207 \\ \hline
	3  & Grande-Bretagne           &2  &41 &38 &45 &124 \\ \hline
	4  & Etats-Unis                &3  &37 &36 &31 &104 \\ \hline
	5  & Comité paralympique Russe &4  &36 &33 &49 &118 \\ \hline
	6  & Pays-Bas                  &5  &25 &17 &17 &59 \\ \hline
	7  & Ukraine                   &6  &24 &47 &27 &98 \\ \hline
	8  & Brésil                    &7  &22 &20 &30 &72 \\ \hline
	9  & Australie                 &8  &21 &29 &30 &80 \\ \hline
	10 & Italie                    &9  &14 &29 &   &69 \\ \hline
	11 & Azerbaïdjan               &10 &14 &1  &4  &19 \\ \hline
	12 & Japon                     &11 &13 &15 &23 &51 \\ \hline
	13 & Allemagne                 &12 &13 &12 &18 &43 \\ \hline
	14 & Iran                      &13 &12 &11 &1  &24 \\ \hline
	15 & France                    &14 &11 &15 &28 &54 \\ \hline
	16 & Espagne                   &15 &9  &15 &12 &36 \\ \hline
\end{tabularx}
\end{center}

\smallskip

\begin{enumerate}
	\item Combien de médailles d’argent l’Australie a-t-elle obtenues ?
	\item Calculer le nombre de médailles de bronze obtenues par l'Italie.
	\item Quelle formule a pu être saisie en \textsf{F2} avant d’être étirée vers le bas ?
	\item Pour chacune des deux affirmations suivantes, indiquer si elle est vraie ou fausse.

	\emph{On rappelle que les réponses doivent être justifiées.}

	\textbf{Affirmation 1:}

	\og 20\,\% des médailles obtenues par l’équipe de France sont en or. \fg{}

	\smallskip

	\textbf{Affirmation 2:}

	\og La médiane du nombre de médailles d'argent obtenues par ces 15 pays est 29. \fg{}

	\item Aux Jeux paralympiques de Rio en 2016, la prime pour une médaille d’or française était de \np{50000} euros. Pour ceux de Tokyo en 2021, cette prime était de \np{65000} euros.

Quel est le pourcentage d’augmentation de cette prime entre 2016 et 2021 ?

\end{enumerate}

\showthesource
\end{document}
